%```latex
\documentclass[12pt,a4paper]{article}

% ------------------------------------------------
% Packages
% ------------------------------------------------

\usepackage[utf8]{inputenc}
\usepackage[T1]{fontenc}
\usepackage[polish]{babel}
\usepackage{lmodern}
\usepackage{geometry}
\usepackage{amsmath, amssymb}
\usepackage{booktabs}
\usepackage{array}
\usepackage{graphicx}
\usepackage{xcolor}
\usepackage{hyperref}
\usepackage{listings}
\usepackage{enumitem}
\usepackage{float}

\usepackage{listings}
\usepackage{listingsutf8}

\lstset{
    inputencoding=utf8,
    extendedchars=true,
    literate=
    {ą}{{\k{a}}}1
    {ć}{{\'c}}1
    {ę}{{\k{e}}}1
    {ł}{{\l}}1
    {ń}{{\'n}}1
    {ó}{{\'o}}1
    {ś}{{\'s}}1
    {ź}{{\'z}}1
    {ż}{{\.z}}1
    {Ą}{{\k{A}}}1
    {Ć}{{\'C}}1
    {Ę}{{\k{E}}}1
    {Ł}{{\L}}1
    {Ń}{{\'N}}1
    {Ó}{{\'O}}1
    {Ś}{{\'S}}1
    {Ź}{{\'Z}}1
    {Ż}{{\.Z}}1
}

\usepackage{amsmath}
\usepackage[normalem]{ulem}
\usepackage{graphicx}

\geometry{
    margin=2.5cm
}

\hypersetup{
    colorlinks=true,
    linkcolor=blue,
    urlcolor=blue,
    citecolor=blue
}

% ------------------------------------------------
% Code formatting
% ------------------------------------------------

\lstset{
    basicstyle=\ttfamily\small,
    breaklines=true,
    frame=single,
    columns=fullflexible,
    keepspaces=true,
    showstringspaces=false
}

% ------------------------------------------------
% Document
% ------------------------------------------------

\title{\textbf{Markdown, Google Colab i GitHub}}
\author{}
\date{}

\begin{document}

\maketitle

\section{Markdown, Colab i GitHub}

Na początek poznamy trzy narzędzia, które rozwiązują trzy różne problemy:

\begin{table}[H]
\centering
\begin{tabular}{ll}
\toprule
\textbf{Narzędzie} & \textbf{Do czego służy?} \\
\midrule
Markdown & Do prostego zapisywania uporządkowanego tekstu, dokumentacji i raportów. \\
Google Colab & Do łączenia tekstu, kodu, obliczeń i wyników w jednym notebooku. \\
GitHub & Do przechowywania projektu, śledzenia zmian i współpracy. \\
\bottomrule
\end{tabular}
\end{table}

Te trzy elementy bardzo dobrze ze sobą współpracują. Podstawą wszędzie jest przygotowanie opisu w kodzie Markdown. Do tego dochodzić będą obliczenia w Colabie, a następnie przechowanie źródła i wyników w repozytorium GitHub.

\section{Markdown}

\subsection{Co to jest Markdown?}

\textbf{Markdown} to prosty język znaczników służący do formatowania zwykłego tekstu.

Najczęściej zapisujemy go w plikach z rozszerzeniem:

\begin{lstlisting}
.md
\end{lstlisting}

Przykłady:

\begin{lstlisting}
README.md
raport.md
notatki.md
\end{lstlisting}

Plik Markdown jest zwykłym plikiem tekstowym. Można go otworzyć praktycznie w każdym edytorze. Specjalne znaki, takie jak \texttt{\#}, \texttt{*}, \texttt{-} czy \texttt{\$}, informują program wyświetlający dokument, jak ma wyglądać tekst.

\subsection{Dlaczego Markdown jest tak istotny?}

Za każdym razem rozmawiając z chatem AI, tak naprawdę rozmawiamy z plikiem Markdown. AI potrafi interpretować składnię Markdown i generować raporty w tym formacie.

Dlatego też, gdy skopiujesz output, ma on całą masę znaczników (np. \texttt{\#}, \texttt{*}, \texttt{-} czy \texttt{\$}), które są potrzebne do poprawnego renderowania.

GitHub (archiwum programistyczne) zawiera opisy również w Markdown. Podobnie Google Colab (środowisko programistyczne Pythona) potrafi renderować Markdown.

Dzięki temu możemy przygotować raport, który będzie wyglądał dobrze zarówno w przeglądarce, jak i w notebooku.

Warto od razu rozróżnić dwie rzeczy:

\begin{itemize}
    \item \textbf{Źródło Markdown} -- tekst, który wpisujemy do pliku.
    \item \textbf{Wyrenderowany dokument} -- wygląd tego tekstu po interpretacji składni Markdown.
\end{itemize}

Na GitHubie możemy przełączać się między widokiem źródła i wyrenderowanym dokumentem.

Podobnie edytory kodu, jak VS Code, PyCharm czy Jupyter Notebook, potrafią renderować Markdown w osobnym oknie podglądu.

\subsection{Przykładowy plik Markdown}

Raport na temat ciała na sprężynie:

\begin{itemize}
    \item \href{examples/01_markdown/raport_cialo_na_sprezynie.md}
    {Ciało na sprężynie}.
\end{itemize}

\subsection{Nagłówki}

Nagłówki budują strukturę dokumentu. Im więcej znaków \texttt{\#}, tym niższy poziom nagłówka.

Źródło:

\begin{lstlisting}
# Tytuł dokumentu

## Główna sekcja

### Podsekcja
\end{lstlisting}

Efekt:

\begin{quote}
    \section*{Tytuł dokumentu}
    \subsection*{Główna sekcja}
    \subsubsection*{Podsekcja}
\end{quote}

W praktyce jeden dokument powinien mieć jeden główny nagłówek \texttt{\#}.

Dalsze części dokumentu organizujemy za pomocą
\texttt{\#\#}, \texttt{\#\#\#} itd.

\subsection{Pogrubienie, kursywa i przekreślenie}

Źródło:

\begin{lstlisting}
**tekst pogrubiony**

*tekst pochylony*

~~tekst przekreślony~~

***tekst pogrubiony i pochylony***
\end{lstlisting}

Efekt:

\begin{itemize}
    \item \textbf{tekst pogrubiony}
    \item \textit{tekst pochylony}
    \item \sout{tekst przekreślony}
    \item \textbf{\textit{tekst pogrubiony i pochylony}}
\end{itemize}

Wyróżnień warto używać oszczędnie. Jeśli wszystko jest pogrubione, nic nie jest naprawdę wyróżnione.

\subsection{Akapit i nowa linia}

Nowy akapit tworzymy przez pozostawienie pustej linii.

Źródło:

\begin{lstlisting}
To jest pierwszy akapit.

To jest drugi akapit.
\end{lstlisting}

Efekt:

To jest pierwszy akapit.

To jest drugi akapit.

To ważna zasada: w Markdown pojedyncze naciśnięcie Enter nie zawsze oznacza nowy akapit.

\subsection{Listy punktowane}

Źródło:

\begin{lstlisting}
- Markdown
- Colab
- GitHub
\end{lstlisting}

Efekt:

\begin{itemize}
    \item Markdown
    \item Colab
    \item GitHub
\end{itemize}

Listy można zagnieżdżać:

\begin{lstlisting}
- GitHub
  - repozytorium
  - Issues
  - historia zmian
- Colab
  - tekst
  - kod
  - wyniki
\end{lstlisting}

Efekt:

\begin{itemize}
    \item GitHub
    \begin{itemize}
        \item repozytorium
        \item Issues
        \item historia zmian
    \end{itemize}
    \item Colab
    \begin{itemize}
        \item tekst
        \item kod
        \item wyniki
    \end{itemize}
\end{itemize}

\subsection{Listy numerowane}

Źródło:

\begin{lstlisting}
1. przygotuj dane,
2. wykonaj obliczenia,
3. zapisz wynik,
4. opisz wynik.
\end{lstlisting}

Efekt:

\begin{enumerate}
    \item przygotuj dane,
    \item wykonaj obliczenia,
    \item zapisz wynik,
    \item opisz wynik.
\end{enumerate}

Listy numerowane są szczególnie przydatne przy opisie procedury lub instrukcji.

\subsection{Checklisty}

GitHub obsługuje listy zadań.

Źródło:

\begin{lstlisting}
- [x] utworzono notebook
- [x] wykonano obliczenia
- [ ] sprawdzono wynik
\end{lstlisting}

Efekt:

\begin{itemize}[label=$\square$]
    \item[\checkmark] utworzono notebook
    \item[\checkmark] wykonano obliczenia
    \item sprawdzono wynik
\end{itemize}

Checklisty bardzo dobrze nadają się do krótkich list kontrolnych w dokumentacji i GitHub Issues.

\subsection{Linki}

Źródło:

\begin{lstlisting}
[GitHub](https://github.com/)
\end{lstlisting}

Efekt:

\href{https://github.com/}{GitHub}

Zamiast wklejać długi adres do strony, warto umieścić go pod czytelną nazwą.

\subsection{Obrazy}

Obraz osadzamy podobnie jak link, ale na początku dodajemy znak \texttt{!}.

Źródło:

\begin{lstlisting}
![Wykres zależności y od x](wykres.png)
\end{lstlisting}

Jeżeli plik \texttt{wykres.png} znajduje się w tym samym folderze co dokument Markdown, GitHub wyświetli obraz w miejscu tej instrukcji.

Jeżeli plik znajduje się w podfolderze \texttt{files}, używamy ścieżki:

\begin{lstlisting}
![Wykres zależności y od x](files/wykres.png)
\end{lstlisting}

Efekt:

\begin{figure}[H]
    \centering
    \includegraphics[width=0.8\textwidth]{wykres.png}
    \caption{Wykres zależności y od x}
\end{figure}

Tekst w nawiasach kwadratowych jest opisem obrazu. Warto go uzupełniać -- pomaga zrozumieć zawartość także wtedy, gdy obraz nie może zostać wyświetlony.

W projektach przechowywanych w repozytorium zwykle wygodniej używać \textbf{ścieżek względnych} niż pełnych adresów internetowych.

\subsection{Tabele}

Źródło:

\begin{lstlisting}
| Narzędzie | Typ pliku | Zastosowanie |
| --- | --- | --- |
| Markdown | .md | dokumentacja |
| Colab | .ipynb | obliczenia |
| GitHub | repozytorium | historia projektu |
\end{lstlisting}

Efekt:

\begin{table}[H]
\centering
\begin{tabular}{lll}
\toprule
\textbf{Narzędzie} & \textbf{Typ pliku} & \textbf{Zastosowanie} \\
\midrule
Markdown & \texttt{.md} & dokumentacja \\
Colab & \texttt{.ipynb} & obliczenia \\
GitHub & repozytorium & historia projektu \\
\bottomrule
\end{tabular}
\end{table}

Tabele są dobre do krótkich zestawień. Jeśli w komórkach pojawiają się długie akapity tekstu, zwykle lepiej wrócić do zwykłych sekcji i list.

\subsection{Kod w tekście}

Nazwy plików, komendy i bardzo krótkie fragmenty kodu warto oznaczać pojedynczym backtickiem.

Źródło:

\begin{lstlisting}
Uruchom plik `analiza.py` i sprawdź wartość zmiennej `wynik`.
\end{lstlisting}

Efekt:

Uruchom plik \texttt{analiza.py} i sprawdź wartość zmiennej \texttt{wynik}.

\subsection{Bloki kodu}

Większy fragment kodu umieszczamy w osobnym bloku.

Źródło:

\begin{lstlisting}
```python
x = 5
y = x**2
print(y)
```
\end{lstlisting}

Efekt:

\begin{lstlisting}[language=Python]
x = 5
y = x**2
print(y)
\end{lstlisting}

Podanie nazwy języka, np. \texttt{python}, \texttt{bash}, \texttt{html} czy \texttt{json}, pozwala GitHubowi kolorować składnię.

W raporcie technicznym dobrze jest pokazać krótki fragment kodu, który jest istotny dla opisywanego wyniku. Nie trzeba kopiować do raportu całego programu.

\subsection{Cytaty}

Źródło:

\begin{lstlisting}
> Wynik należy zawsze sprawdzić przed jego wykorzystaniem.
\end{lstlisting}

Efekt:

\begin{quote}
Wynik należy zawsze sprawdzić przed jego wykorzystaniem.
\end{quote}

Cytaty można wykorzystać także do krótkiego wyróżnienia komentarza lub uwagi.

\subsection{Matematyka}

GitHub potrafi renderować wzory matematyczne.

Krótki wzór umieszczamy wewnątrz pojedynczych znaków dolara.

Źródło:

\begin{lstlisting}
Energia spoczynkowa dana jest wzorem $E=mc^2$.
\end{lstlisting}

Efekt:

Energia spoczynkowa dana jest wzorem \(E=mc^2\).

Dłuższy wzór umieszczamy w osobnym bloku.

Źródło:

\begin{lstlisting}
$$
\sum_{k=1}^{n} k^2 = \frac{n(n+1)(2n+1)}{6}
$$
\end{lstlisting}

Efekt:

\[
\sum_{k=1}^{n} k^2 =
\frac{n(n+1)(2n+1)}{6}
\]

Macierze zapisujemy wielowierszowo:

Źródło:

\begin{lstlisting}
$$
A=
\begin{pmatrix}
1 & 2 \\
3 & 4
\end{pmatrix}
$$
\end{lstlisting}

Efekt:

\[
A=
\begin{pmatrix}
1 & 2 \\
3 & 4
\end{pmatrix}
\]

\subsection{Zrzut ekranu i transkrypcja do Markdown}

Czasem informacja istnieje tylko jako obraz: zrzut ekranu programu, zdjęcie tabeli, komunikat błędu albo fragment zeskanowanego dokumentu.

AI może pomóc przepisać zawartość obrazu do tekstu. Taki proces nazywamy tutaj \textbf{transkrypcją}.

Przykładowy workflow:

\begin{enumerate}
    \item Zachowujemy oryginalny obraz.
    \item Prosimy AI o \textbf{przepisanie jego zawartości do kodu Markdown}.
    \item Zapisujemy wynik w pliku tekstowym lub Markdown.
    \item Porównujemy transkrypcję z obrazem.
    \item Poprawiamy ewentualne błędy.
\end{enumerate}

Największej ostrożności wymagają:

\begin{itemize}
    \item liczby,
    \item wzory,
    \item nazwy plików,
    \item adresy,
    \item fragmenty kodu,
    \item komunikaty błędów.
\end{itemize}

Transkrypcja wygenerowana przez AI jest \textbf{wersją roboczą}, dopóki człowiek jej nie sprawdzi.

\subsection{Poprawne formatowanie matematyki w Markdown}

W materiałach tego kursu polecamy stosowanie prostego zestawu reguł formatowania matematyki przygotowanego specjalnie tak, aby pliki dobrze renderowały się na GitHubie:

\begin{itemize}
    \item \href{hints.md}{hints.md}
\end{itemize}

Pracując z chatem AI i domagając się kodu Markdown, warto załączać ten plik jako kontekst, aby AI od razu generowało poprawny kod w Markdown.

Różne systemy: przeglądarkowe chaty AI, edytor Visual Studio Code, Colab i GitHub mogą różnie renderować wzory matematyczne!

Częstym błędem jest używanie w kodzie Markdown znaków \texttt{\textbackslash[} i \texttt{\textbackslash]} zamiast \texttt{\$...\$} do oznaczania wzorów matematycznych. Wtedy GitHub nie renderuje wzoru, a Colab renderuje go poprawnie.

\subsection{Dobre materiały o Markdown}

\begin{itemize}
    \item \href{https://docs.github.com/en/get-started/writing-on-github/getting-started-with-writing-and-formatting-on-github/basic-writing-and-formatting-syntax}
    {GitHub Docs -- Basic writing and formatting syntax}

    \item \href{https://docs.github.com/en/get-started/writing-on-github/working-with-advanced-formatting/writing-mathematical-expressions}
    {GitHub Docs -- Writing mathematical expressions}

    \item \href{https://www.markdownguide.org/basic-syntax/}
    {Markdown Guide -- Basic Syntax}

    \item \href{https://www.markdownguide.org/cheat-sheet/}
    {Markdown Guide -- Cheat Sheet}
\end{itemize}



\end{document}

\end{document}
```